\chapter{Compiling this document and regenerating the figures}\label{app:build}

\section{Compiling the PDF}
The report uses XeLaTeX-compatible Unicode fonts and the \texttt{listings} package. Code is rendered in a continuous light-blue panel with line numbers; the number column is italic and grey. No shell escape or network access is required.

\begin{lstlisting}[style=vstyle,language=text,numbers=left,firstnumber=1]
cd documentation/documentation_sources/latex
tectonic main.tex
# or, with TeX Live:
latexmk -xelatex main.tex
\end{lstlisting}

Run the compiler until the table of contents, list of figures and cross-references stabilise. The source tree is:

\begin{lstlisting}[style=vstyle,language=text,numbers=left,firstnumber=1]
main.tex  ch01_intro.tex ... ch19_missy.tex  chV1_method.tex ... chV5_findings.tex
appA_build.tex ... appI_atlas_data.tex
fig/  screenshots and generated figures
code/ complete HTML and verification extracts
tables/ generated tables
data/atlas/ CSV values behind the control-chart atlas
\end{lstlisting}

\section{Regenerating figures and tables}
The generator uses Node.js, Python and Playwright. The supplied verification script loads the production HTML in a real browser, captures the three modes and checks the report exports.

\begin{lstlisting}[style=vstyle,language=text,numbers=left,firstnumber=1]
node documentation/verification/missy_smoke_2026-09-25.cjs
\end{lstlisting}

The atlas source is \path{DEMO_control_charts_realscale.zip}. Appendix~\ref{app:atlas-data} lists every CSV and records the manifest without reproducing CSV listings. The companion A4 PDF contains one chart per page.



